\chapter{Classification of Differential Equations}

% -------------------------------------------------------
\section{Order}
% -------------------------------------------------------

\begin{definition}[Order]
The \textbf{order} of a differential equation is the order of the highest derivative present.
\end{definition}

\begin{example}
\begin{align*}
  \frac{dy}{dx} + y = 0 &\quad \text{(first order)} \\[6pt]
  \frac{d^2y}{dx^2} + 3\frac{dy}{dx} + y = x &\quad \text{(second order)} \\[6pt]
  \frac{d^3y}{dx^3} = \sin x &\quad \text{(third order)}
\end{align*}
\end{example}

% -------------------------------------------------------
\section{Linearity}
% -------------------------------------------------------

\subsection{Linear Differential Equation}

\begin{definition}[Linear ODE]
An ODE is \textbf{linear} if:
\begin{enumerate}
  \item the dependent variable and all its derivatives appear to the first power, and
  \item there is no multiplication between the dependent variable and any of its derivatives.
\end{enumerate}
The general form of a linear $n$-th order ODE is
\[
  a_n(x)\,y^{(n)} + a_{n-1}(x)\,y^{(n-1)} + \cdots + a_1(x)\,y' + a_0(x)\,y = f(x).
\]
\end{definition}

\begin{example}
\[
  \frac{dy}{dx} + 3y = x
\]
This is linear: $y$ and $dy/dx$ each appear to the first power and are not multiplied together.
\end{example}

\subsection{Non-Linear Differential Equation}

\begin{definition}[Non-Linear ODE]
An ODE is \textbf{non-linear} if the dependent variable or its derivatives appear in higher powers, under a non-linear function, or are multiplied together.
\end{definition}

\begin{example}
\[
  y\,\frac{dy}{dx} = x
\]
This is non-linear because $y$ is multiplied by $dy/dx$.
\end{example}

% -------------------------------------------------------
\section{Homogeneous vs.\ Non-Homogeneous}
% -------------------------------------------------------

\begin{definition}[Homogeneous Equation]
A linear ODE is \textbf{homogeneous} if the right-hand side is identically zero:
\[
  a_n(x)\,y^{(n)} + \cdots + a_0(x)\,y = 0.
\]
\end{definition}

\begin{definition}[Non-Homogeneous Equation]
A linear ODE is \textbf{non-homogeneous} if the right-hand side is a non-zero function:
\[
  a_n(x)\,y^{(n)} + \cdots + a_0(x)\,y = f(x), \qquad f(x) \not\equiv 0.
\]
\end{definition}

\begin{example}
\begin{align*}
  y'' + y = 0 &\quad \text{(homogeneous)}\\
  y'' + y = \sin x &\quad \text{(non-homogeneous, } f(x)=\sin x \text{)}
\end{align*}
\end{example}
